% =============================================================================
% Chapter 6: Network Architectures — ML Systems Lecture Slides
% =============================================================================
% IMAGES:
%   - ch06-inductive-bias-spectrum.pdf
%   - ch06-architecture-bottleneck-map.pdf
%   - ch06-cnn-weight-sharing.pdf
%   - ch06-attention-mechanism.pdf
%   - ch06-transformer-block.pdf
%   - ch06-rnn-unrolled.pdf
%   - ch06-architecture-selection-flowchart.pdf
% =============================================================================
\documentclass[aspectratio=169, 12pt]{beamer}
\usepackage{../../assets/beamerthememlsys}

\mlsyssetup{
  volume   = {Volume I},
  chapter  = {Chapter 6},
  logo = {../../assets/img/logo-mlsysbook.png},
  instlogo = {../../assets/img/logo-harvard.png},
  chaptertitle = {Network Architectures},
}

% --- Fonts ---
\usepackage{fontspec}
\setsansfont{Helvetica Neue}[
  BoldFont={Helvetica Neue Bold},
  ItalicFont={Helvetica Neue Italic},
  BoldItalicFont={Helvetica Neue Bold Italic},
]
\setmonofont{JetBrains Mono}[Scale=0.85]

% --- Packages ---
\usepackage{booktabs}
\usepackage{amsmath}

% --- Image paths ---
\graphicspath{
  {images/}
}

% --- Chapter-specific macros ---
\newcommand{\DAM}{D\raisebox{0.08em}{\tiny$\bullet$}A\raisebox{0.08em}{\tiny$\bullet$}M}

% --- Helper: safe image include ---
\newcommand{\safeimg}[2][width=\textwidth,keepaspectratio]{%
  \IfFileExists{images/#2}{\includegraphics[#1]{#2}}{%
    \IfFileExists{#2}{\includegraphics[#1]{#2}}{%
      \fbox{\parbox[c][2.5cm][c]{0.85\linewidth}{\centering\footnotesize\textcolor{midgray}{[Missing image]}}}%
    }%
  }%
}

% --- Section count for navigation ---
\setcounter{mlsystotalsections}{7}

\title{Network Architectures}
\author{Vijay Janapa Reddi}
\institute{Harvard University}
\date{}

\begin{document}

% =============================================================================
% TITLE SLIDE
% =============================================================================
\mlsystitle{Network Architectures}{Architecture Is Infrastructure}{cover_dl_arch.png}

% =============================================================================
% LEARNING OBJECTIVES
% =============================================================================
\begin{frame}{Learning Objectives}
\note{
% -- NARRATE: What to SAY while showing this slide
Walk through each objective. Stress that this chapter is about the
\emph{systems consequences} of architecture choice, not model accuracy.
``Every architecture signs a contract with physics---by the end of today,
you will be able to read that contract.''

% -- ENGAGE: Specific question, prediction, or task for THIS slide
Ask: ``When you pick a Transformer, what hardware contract are you signing?''
Cold-call one student. [Expected answer: O(N\^{}2) memory, high bandwidth demand.]

% -- FLEX: [CORE] This slide is essential --- do not skip.
[CORE] Sets the framing for the entire lecture.
IF SHORT: Read objectives 1, 4, and 6 aloud; let students read the rest.
}

\small
\begin{enumerate}\setlength\itemsep{0pt}
  \item Distinguish the \textbf{computational signatures} of MLPs, CNNs, RNNs, Transformers, and DLRM
  \item Explain how \textbf{inductive biases} encode data structure into computation
  \item Analyze \textbf{arithmetic intensity} and memory scaling across architectures
  \item Identify shared \textbf{building blocks} (skip connections, normalization, gating)
  \item Apply the \textbf{architecture selection framework} to match data, hardware, and constraints
  \item Evaluate how architecture determines \textbf{Iron Law} bottlenecks ($O$, $D_{\text{vol}}$, $L_{\text{lat}}$)
  \item Critique common architectural selection \textbf{fallacies} using systems reasoning
\end{enumerate}

\end{frame}

\begin{frame}{Visual Language}
\note{
% -- NARRATE: What to SAY while showing this slide
Point to each card: ``Blue means compute---anytime you see blue in a diagram,
that is where the GPU is doing math. Green is data flow. Orange is scheduling
decisions. Red flags a bottleneck or cost.''

% -- FLEX: [CORE] Establishes the visual vocabulary for every subsequent slide.
[CORE] First time students see the color system; revisit briefly at start of Ch7.
IF SHORT: Show for 30 seconds and move on---students will internalize through use.
}

\small
Throughout this course, colors carry meaning:

\vspace{0.3cm}
\begin{columns}[T]
  \begin{column}{0.45\textwidth}
    \begin{mlsyscard}{computestroke}
    \textbf{Blue} --- Compute / Processing\\
    {\footnotesize GPU ops, forward/backward pass, inference}
    \end{mlsyscard}
    \vspace{0.15cm}
    \begin{mlsyscard}{datastroke}
    \textbf{Green} --- Data / Memory\\
    {\footnotesize Data flow, caches, healthy paths}
    \end{mlsyscard}
  \end{column}
  \begin{column}{0.45\textwidth}
    \begin{mlsyscard}{routingstroke}
    \textbf{Orange} --- Routing / Scheduling\\
    {\footnotesize Load balancers, batch windows}
    \end{mlsyscard}
    \vspace{0.15cm}
    \begin{mlsyscard}{errorstroke}
    \textbf{Red} --- Error / Cost / Bottleneck\\
    {\footnotesize Loss, decode phase, waste}
    \end{mlsyscard}
  \end{column}
\end{columns}
\end{frame}


% =============================================================================
\section{Inductive Bias}
% =============================================================================

\begin{frame}{Architecture Is a Contract with Physics}
\note{
% -- LINK: What prior concept connects to this slide
Students learned the Iron Law ($O$, $D_{\text{vol}}$, $L_{\text{lat}}$) and the
roofline model in earlier chapters. This slide shows that the architecture
\emph{determines} which Iron Law term dominates.

% -- NARRATE: What to SAY while showing this slide
``Architecture is not a modeling decision---it is an infrastructure commitment.
When you pick a CNN, you commit to spatial locality and high arithmetic intensity.
When you pick a Transformer, you commit to O(N\^{}2) memory.''
Point to the table on the right: ``RNN: 30--50\% GPU utilization. That is the
contract you sign.''

% -- ENGAGE: Specific question, prediction, or task for THIS slide
Ask: ``Has anyone picked an architecture and later regretted the systems cost?''
Let 2 students share experiences before continuing.

% -- WARN: What students will get wrong on THIS topic
Common error: students choose architectures based on accuracy benchmarks alone,
ignoring memory, latency, and power constraints.
Correct framing: architecture determines the physical resources required to
run the model---accuracy is necessary but not sufficient.

% -- FLEX: [CORE] This slide is the chapter thesis.
[CORE] Do not skip---every subsequent slide builds on this framing.
IF AHEAD: ``Can you name an architecture whose systems cost killed a product?''
}

\footnotesize
\begin{columns}[T]
  \begin{column}{0.55\textwidth}
    Choosing an architecture means \alert{signing a contract}:

    \vspace{0.1cm}
    \begin{itemize}\setlength\itemsep{0pt}
      \item \textbf{CNN} $\to$ spatial locality, parallel convolutions
      \item \textbf{Transformer} $\to$ $O(N^2)$ memory, global attention
      \item \textbf{RNN} $\to$ sequential dependencies, low parallelism
      \item \textbf{DLRM} $\to$ TB-scale embedding tables
    \end{itemize}

    \vspace{0.1cm}
    \begin{mlsyscard}{crimson}
    Architecture determines \textbf{Iron Law terms}:
    Operations ($O$), data movement ($D_{\text{vol}}$), overhead ($L_{\text{lat}}$).
    \end{mlsyscard}
  \end{column}
  \begin{column}{0.42\textwidth}
    \renewcommand{\arraystretch}{1.15}
    {\footnotesize
    \begin{tabular}{@{}ll@{}}
      \toprule
      \textbf{Choice} & \textbf{Consequence} \\
      \midrule
      CNN    & Fits mobile memory \\
      Transf.& Needs datacenter \\
      RNN    & 30--50\% GPU util. \\
      DLRM   & TB of DRAM \\
      \bottomrule
    \end{tabular}
    }
  \end{column}
\end{columns}

\end{frame}

% --- ACTIVE LEARNING 1: Predict ---
\begin{frame}{Predict: What Unifies All Architectures?}
\note{
% -- LINK: What prior concept connects to this slide
The previous slide showed four architectures with wildly different contracts.
Now students must identify what single concept explains \emph{why} those
contracts differ.

% -- NARRATE: What to SAY while showing this slide
``Four architectures, four very different hardware contracts. There is one
concept that explains all of them. Write your answer---60 seconds.''
Do NOT reveal ``inductive bias'' yet.

% -- ENGAGE: Specific question, prediction, or task for THIS slide
Students write individually for 60 seconds. Then share with a neighbor.
Cold-call 2--3 students to share answers before moving to the reveal slide.
[Expected answer: inductive bias---the structural assumption each architecture
encodes about data.]

% -- FLEX: [CORE] This prediction primes the key concept reveal.
[CORE] The prediction must happen before the reveal slide.
IF SHORT: Reduce to 30 seconds writing, skip neighbor sharing.
}

\centering
\vspace{1.0cm}
{\Large\bfseries Think--Write--Share}

\vspace{0.6cm}
{\large MLPs, CNNs, RNNs, and Transformers look very different.\\
What single concept \emph{unifies} all architecture choices?}

\vspace{0.5cm}
{\normalsize Write your answer. \textcolor{midgray}{(60 seconds)}}

\vspace{0.5cm}
{\small\textcolor{lightgray}{Hint: it explains why CNNs use 111$\times$ fewer parameters than MLPs.}}

\end{frame}

\begin{frame}{Inductive Bias: The Unifying Concept}
\note{
% -- LINK: What prior concept connects to this slide
Students just predicted the unifying concept. This slide reveals the answer:
inductive bias. Connect to the hint---111x fewer parameters is the consequence
of the locality bias in CNNs.

% -- NARRATE: What to SAY while showing this slide
Point to the spectrum diagram left to right: ``Stronger bias means fewer
parameters means less flexibility but better hardware mapping. An MLP has
no bias---it must learn everything from scratch. A CNN assumes spatial
locality and saves 111x parameters. The bias determines which Iron Law
term dominates.''
ANALOGY: ``Inductive bias is like building codes---they constrain what you
can build, but the constraints make the building possible.''

% -- ENGAGE: Specific question, prediction, or task for THIS slide
Ask: ``Where on this spectrum would you place a Graph Neural Network?''
Give 15 seconds. [Expected answer: between CNN and Transformer---structural
bias on graph topology, but more flexible than spatial grids.]

% -- WARN: What students will get wrong on THIS topic
Common error: students think stronger bias is always better.
Correct framing: stronger bias reduces parameters but restricts what the
model can learn. Transformers sacrifice efficiency for flexibility.
IF STUCK: ``Think of bias as a constraint---when does a constraint help
vs.\ hurt?''

% -- FLEX: [CORE] This is the chapter's central definition.
[CORE] Every subsequent architecture slide references inductive bias.
IF AHEAD: ``What happens when your bias mismatches the data structure?''
}

% --- Layout: FULL-WIDTH diagram ---
\centering
\safeimg[width=0.92\textwidth,height=4.8cm,keepaspectratio]{ch06-inductive-bias-spectrum.pdf}

\vspace{0.15cm}
\mlsysconcept{Definition:}{Inductive bias is a structural constraint that restricts the hypothesis space, encoding domain assumptions directly into the computational graph.}

\end{frame}

% =============================================================================
\section{Architecture Zoo}
% =============================================================================

\begin{frame}{Five Lighthouse Models}
\note{
% -- LINK: What prior concept connects to this slide
Inductive bias explains \emph{why} architectures differ. Now we introduce
five concrete architectures that each isolate a different system bottleneck,
mapping to the workload archetypes from Ch2.

% -- NARRATE: What to SAY while showing this slide
Point to each lighthouse in the diagram: ``ResNet isolates compute.
GPT-2 isolates bandwidth. DLRM isolates memory capacity. MobileNet
isolates latency on edge. KWS isolates power on microcontrollers.
No single architecture is best---each teaches us about a different
physical constraint.''

% -- ENGAGE: Specific question, prediction, or task for THIS slide
Ask: ``Why do we need five lighthouse models, not one?''
Cold-call. [Expected answer: each isolates a different bottleneck;
one model cannot teach all system constraints.]

% -- WARN: What students will get wrong on THIS topic
Common error: students think ``lighthouse'' means ``best'' architecture.
Correct framing: lighthouses are diagnostic tools, not recommendations.
Each reveals a specific hardware-workload interaction.

% -- FLEX: [CORE] Establishes the five recurring examples for the rest of the course.
[CORE] Students must memorize which bottleneck each lighthouse isolates.
IF SHORT: Name all five quickly; spend detail time on the next slides.
}

% --- Layout: FULL-WIDTH diagram ---
\centering
\safeimg[width=0.92\textwidth,height=4.8cm,keepaspectratio]{ch06-architecture-bottleneck-map.pdf}

\vspace{0.1cm}
{\small Each lighthouse isolates a \textbf{distinct system bottleneck} --- no single ``best'' architecture exists.}

\end{frame}

\begin{frame}{Arithmetic Intensity Spectrum}
\note{
% -- LINK: What prior concept connects to this slide
The five lighthouses each live at different points on the roofline model
from Ch1. Arithmetic intensity (FLOPs/byte) determines whether a workload
is compute-bound or memory-bound.

% -- NARRATE: What to SAY while showing this slide
Point to each row: ``ResNet at 40 FLOPs/byte---above the ridge point,
compute-bound, GPU-friendly. GPT-2 at 0.5 FLOPs/byte---below the ridge,
bandwidth-starved. Three orders of magnitude separate them. The architecture
determines whether your system needs a faster processor or faster memory.''

% -- ENGAGE: Specific question, prediction, or task for THIS slide
Ask: ``If MobileNet has 21 FLOPs/byte and a datacenter GPU ridge point is
156 FLOPs/byte, what regime is MobileNet in on that GPU?''
Give 10 seconds. [Expected answer: memory-bound---it starves the GPU.]

% -- WARN: What students will get wrong on THIS topic
Common error: students assume high FLOPs = fast execution.
Correct framing: speed depends on arithmetic intensity relative to the
hardware's ridge point, not absolute FLOP count.
IF STUCK: ``Draw the roofline---where does each model land?''

% -- FLEX: [CORE] Arithmetic intensity is the quantitative bridge between
architecture and hardware.
[CORE] This table is referenced in the Fallacies slide and the selection framework.
IF AHEAD: ``What would it take to make GPT-2 inference compute-bound?''
}

\small
\textbf{Arithmetic Intensity} ($I$) = FLOPs per byte of data moved

\vspace{0.15cm}
\renewcommand{\arraystretch}{1.15}
{\footnotesize
\begin{tabular}{@{}llll@{}}
  \toprule
  \textbf{Model} & \textbf{Intensity} & \textbf{Regime} & \textbf{Hardware Affinity} \\
  \midrule
  \textbf{ResNet-50} & $\sim$40 FLOPs/B & \textcolor{computestroke}{Compute-bound} & GPU/TPU (saturates ALUs) \\
  \textbf{MobileNet} & $\sim$21 FLOPs/B & Balanced & Mobile NPU \\
  \textbf{GPT-2 (inf)} & $\sim$0.5 FLOPs/B & \textcolor{datastroke}{Memory-bound} & HBM3 bandwidth \\
  \textbf{DLRM} & $\ll 1$ FLOPs/B & \textcolor{errorstroke}{Capacity-bound} & TB-scale DRAM \\
  \bottomrule
\end{tabular}
}

\vspace{0.2cm}
\begin{mlsyscard}{crimson}
{\footnotesize \textbf{Three orders of magnitude} separate ResNet and GPT-2. The architecture determines whether your system needs a faster processor or faster memory.}
\end{mlsyscard}

\end{frame}

\begin{frame}{MLPs: The Universal Baseline}
\note{
% -- LINK: What prior concept connects to this slide
The inductive bias spectrum placed MLPs at the far left---zero bias, full
connectivity. This slide unpacks what that means in systems terms.

% -- NARRATE: What to SAY while showing this slide
``MLPs are the simplest architecture: no structural assumption, every input
connects to every output. The Universal Approximation Theorem guarantees
they can represent anything, but at enormous cost. One 1024-to-1024 layer:
over one million parameters. That is the price of zero inductive bias.''
Point to the table: ``High parallelism and 80--90\% utilization---the
computation is simple dense MatMul, which GPUs love.''

% -- ENGAGE: Specific question, prediction, or task for THIS slide
Ask: ``If MLPs can approximate any function, why do we need CNNs or
Transformers at all?''
Give 15 seconds. [Expected answer: MLPs require exponentially more parameters
to learn structure that biased architectures encode for free.]

% -- WARN: What students will get wrong on THIS topic
Common error: students dismiss MLPs as outdated. In practice, MLP layers
are the computational workhorse inside every Transformer (the FFN block).
Correct framing: MLPs are the universal building block; other architectures
add constraints on top.

% -- FLEX: [CORE] Establishes the baseline against which all other architectures
are compared.
[CORE] The 1M parameter number recurs in the CNN comparison.
IF SHORT: Skip the UAT card; keep the table and the ``baseline'' insight.
}

\small
\begin{columns}[T]
  \begin{column}{0.55\textwidth}
    \textbf{Multi-Layer Perceptrons}
    \begin{itemize}\setlength\itemsep{0pt}
      \item {\footnotesize \textbf{Bias}: None --- any input relates to any output}
      \item {\footnotesize \textbf{Data}: Tabular, unstructured features}
      \item {\footnotesize \textbf{Scaling}: $O(d^2)$ params per layer}
      \item {\footnotesize \textbf{1024 $\to$ 1024}: 1,048,576 parameters}
    \end{itemize}

    \vspace{0.1cm}
    \begin{mlsyscard}{errorstroke}
    {\footnotesize UAT guarantees MLPs \emph{can} represent any function. Catch: ``sufficiently wide'' can mean exponentially many neurons.}
    \end{mlsyscard}
  \end{column}
  \begin{column}{0.42\textwidth}
    \renewcommand{\arraystretch}{1.05}
    {\scriptsize
    \begin{tabular}{@{}lr@{}}
      \toprule
      \textbf{Property} & \textbf{Value} \\
      \midrule
      Bias & None \\
      Params/layer & $O(d^2)$ \\
      Parallelism & High \\
      Utilization & 80--90\% \\
      Bottleneck & Bandwidth \\
      \bottomrule
    \end{tabular}
    }

    \vspace{0.1cm}
    \mlsysinsight{Baseline:}{Every other architecture is an MLP + structural constraints.}
  \end{column}
\end{columns}

\end{frame}

\begin{frame}{CNNs: Spatial Pattern Processing}
\note{
% -- LINK: What prior concept connects to this slide
MLPs required 1M parameters for a 1024x1024 layer. CNNs apply the locality
bias to achieve the same task with 111x fewer parameters.

% -- NARRATE: What to SAY while showing this slide
Walk through the weight sharing diagram: ``Left side: MLP connects every
pixel to every output---1,048,576 parameters. Right side: a single 3x3
filter slides across the entire image---9,408 parameters. Same spatial
task, 111x cheaper. That is the power of the locality bias: if nearby
pixels matter most, do not waste parameters on distant connections.''

% -- ENGAGE: Specific question, prediction, or task for THIS slide
Ask: ``What happens to CNN accuracy if you shuffle the pixel order of
every image?'' Give 10 seconds.
[Expected answer: accuracy collapses---the locality bias assumes spatial
structure. Shuffled pixels destroy that structure.]

% -- WARN: What students will get wrong on THIS topic
Common error: students think weight sharing only saves parameters.
Correct framing: weight sharing also increases arithmetic intensity
because each weight is reused across many spatial positions, making
CNNs compute-bound and GPU-friendly.

% -- FLEX: [CORE] The 111x number is a key quantitative anchor.
[CORE] Referenced in Key Takeaways and the Fallacies slide.
IF AHEAD: ``How does depthwise separable convolution change this ratio?''
}

% --- Layout: FULL-WIDTH diagram ---
\centering
\safeimg[width=0.92\textwidth,height=4.5cm,keepaspectratio]{ch06-cnn-weight-sharing.pdf}

\vspace{0.1cm}
\small
\begin{columns}[T]
  \begin{column}{0.48\textwidth}
    \centering
    \textcolor{errorstroke}{\textbf{MLP}}: 1,048,576 parameters\\
    {\footnotesize Every pixel $\to$ every output}
  \end{column}
  \begin{column}{0.48\textwidth}
    \centering
    \textcolor{datastroke}{\textbf{CNN}}: 9,408 parameters\\
    {\footnotesize Shared 3$\times$3 filter slides across image}
  \end{column}
\end{columns}
\vspace{0.05cm}
\mlsysconcept{Key:}{Weight sharing = 111$\times$ fewer params = fits on mobile devices.}

\end{frame}

\begin{frame}{CNN System Properties}
\note{
% -- LINK: What prior concept connects to this slide
Weight sharing (previous slide) explains why CNNs have fewer parameters.
This slide shows the systems consequence: high arithmetic intensity and
excellent GPU utilization.

% -- NARRATE: What to SAY while showing this slide
``Because each weight is reused many times across spatial positions, CNNs
achieve high arithmetic intensity---around 40 FLOPs/byte for ResNet-50.
That puts them above the ridge point: compute-bound, the ideal regime for
GPUs. ResNet-50 achieves 80--95\% of peak throughput.''
Point to the MobileNet column: ``Depthwise separable convolutions cut
FLOPs by 8--9x, making MobileNet fit on phones at 14 MB and 300 MFLOPs.
But the tradeoff: low intensity starves datacenter GPUs.''

% -- ENGAGE: Specific question, prediction, or task for THIS slide
Ask: ``MobileNet uses 14x fewer FLOPs than ResNet. Will it run 14x faster
on an A100?'' Give 10 seconds.
[Expected answer: no---low arithmetic intensity means MobileNet is
memory-bound on datacenter GPUs and may actually run slower.]

% -- WARN: What students will get wrong on THIS topic
Common error: students equate fewer FLOPs with faster execution.
Correct framing: FLOPs predict speed only when compute-bound. MobileNet's
low intensity means the GPU waits for memory, not math.
IF STUCK: ``Draw both models on the roofline---which hits the ceiling?''

% -- FLEX: [CORE] The FLOPs-vs-speed paradox is a central lesson.
[CORE] This directly supports the ``FLOPs != speed'' fallacy slide.
IF SHORT: Skip MobileNet details; keep the ResNet utilization numbers.
}

\footnotesize
\begin{columns}[T]
  \begin{column}{0.52\textwidth}
    \textbf{Why CNNs map well to GPUs:}
    \begin{itemize}\setlength\itemsep{0pt}
      \item Weight reuse $\to$ high arithmetic intensity
      \item Regular memory access $\to$ cache-friendly
      \item Batch processing $\to$ saturates tensor cores
      \item \textbf{ResNet-50}: 80--95\% peak throughput
    \end{itemize}

    \vspace{0.1cm}
    \textbf{MobileNet (Depthwise Separable):}
    \begin{itemize}\setlength\itemsep{0pt}
      \item 8--9$\times$ fewer FLOPs for 3$\times$3 kernels
      \item Fits phones: 14 MB, 300 MFLOPs
      \item But: \alert{starves datacenter GPUs} (low intensity)
    \end{itemize}
  \end{column}
  \begin{column}{0.45\textwidth}
    \renewcommand{\arraystretch}{1.1}
    {\scriptsize
    \begin{tabular}{@{}lrr@{}}
      \toprule
      & \textbf{ResNet-50} & \textbf{MobileNet} \\
      \midrule
      Params     & 25.6 M & 3.5 M \\
      FLOPs      & 4.1 G  & 300 M \\
      Memory     & 98 MB  & 14 MB \\
      Intensity  & $\sim$40  & $\sim$21 \\
      GPU util.  & 80--95\% & 30--50\% \\
      \bottomrule
    \end{tabular}
    }

    \vspace{0.1cm}
    \begin{mlsyscard}{errorstroke}
    {\scriptsize \textbf{FLOPs $\neq$ speed!} MobileNet uses 14$\times$ fewer FLOPs but can run \emph{slower} on datacenter GPUs.}
    \end{mlsyscard}
  \end{column}
\end{columns}

\end{frame}

\begin{frame}{Quick Check}
\note{
% -- NARRATE: What to SAY while showing this slide
``Quick check---no peeking. What is the memory scaling of Transformer
attention with respect to sequence length? 30 seconds.''
After 30 seconds, take one answer aloud. [Expected answer: O(N\^{}2).]

% -- FLEX: [OPTIONAL] Micro-retrieval for spaced practice.
[OPTIONAL] If running behind, skip this and move directly to RNNs.
}

\centering
\Large\bfseries Quick check --- no peeking.\\[0.5cm]
\normalsize What is the memory scaling of Transformer attention with respect to sequence length?\\[0.3cm]
{\small 30 seconds --- then we continue.}
\end{frame}



\begin{frame}{RNNs: Sequential Pattern Processing}
\note{
% -- LINK: What prior concept connects to this slide
CNNs exploited spatial locality for parallelism. RNNs handle sequential
data but pay a steep systems cost: each step depends on the previous one,
killing GPU parallelism.

% -- NARRATE: What to SAY while showing this slide
Walk through the unrolled RNN diagram: ``Each timestep reads the hidden
state from the previous step. The GPU cannot start step $t$ until step
$t-1$ finishes. This sequential dependency limits utilization to 30--50\%.''
Point to the dashed red line: ``The vanishing gradient---by the time
gradients reach early timesteps, they have shrunk exponentially. LSTM
and GRU gates partially solve this but cannot fix the parallelism problem.''
ANALOGY: ``An RNN is like a single-lane highway---no matter how fast each
car goes, throughput is limited by sequential flow.''

% -- ENGAGE: Specific question, prediction, or task for THIS slide
Ask: ``If RNNs have O(1) memory per step, why did Transformers replace
them for language?'' Give 15 seconds.
[Expected answer: O(1) memory is great for streaming, but sequential
processing is too slow for training on modern GPUs.]

% -- WARN: What students will get wrong on THIS topic
Common error: students assume RNNs are obsolete.
Correct framing: RNNs are ideal for streaming edge deployment (O(1) memory,
low latency per step). They fail at training throughput, not at all tasks.

% -- FLEX: [CORE] Sets up the RNN-vs-Transformer trade-off on the next slide.
[CORE] The 30--50\% utilization number is referenced in the comparison table.
IF SHORT: Focus on the sequential bottleneck; skip the vanishing gradient detail.
}

% --- Layout: FULL-WIDTH diagram ---
\centering
\safeimg[width=0.92\textwidth,height=4.5cm,keepaspectratio]{ch06-rnn-unrolled.pdf}

\vspace{0.1cm}
\small
\begin{columns}[T]
  \begin{column}{0.48\textwidth}
    \centering
    \textcolor{routingstroke}{\textbf{Strength}}: $O(1)$ memory per step\\
    {\footnotesize Constant state size, streaming-friendly}
  \end{column}
  \begin{column}{0.48\textwidth}
    \centering
    \textcolor{errorstroke}{\textbf{Weakness}}: Sequential bottleneck\\
    {\footnotesize 30--50\% GPU utilization, gradient vanishes}
  \end{column}
\end{columns}

\end{frame}

% =============================================================================
\section{Attention}
% =============================================================================

\begin{frame}{Attention: Dynamic Processing}
\note{
% -- LINK: What prior concept connects to this slide
RNNs process sequentially with O(1) memory. Attention replaces fixed
sequential processing with learned, content-dependent routing---but at
O(N\^{}2) memory cost.

% -- NARRATE: What to SAY while showing this slide
Walk through the QKV diagram: ``Query asks `what am I looking for?'
Key says `what do I contain?' Value says `here is my content.'
The attention matrix computes relevance between ALL pairs of elements.
For 4096 tokens, that matrix is 4096x4096 = 16M entries. At 2 bytes
each (FP16), that is 32 MB per layer per head.''

% -- ENGAGE: Specific question, prediction, or task for THIS slide
Ask: ``If you double the sequence length from 4096 to 8192, how much
more attention memory do you need?'' Give 10 seconds.
[Expected answer: 4x---quadratic scaling.]

% -- WARN: What students will get wrong on THIS topic
Common error: students think attention is expensive because of the
dot products. In reality, the bottleneck is \emph{materializing} the
NxN matrix in HBM---it is a memory problem, not a compute problem.
IF STUCK: ``Think about where the NxN matrix lives during computation.''

% -- FLEX: [CORE] O(N\^{}2) scaling is the defining constraint of Transformers.
[CORE] This sets up the memory budget exercise two slides later.
IF AHEAD: ``How does FlashAttention avoid materializing this matrix?''
}

% --- Layout: FULL-WIDTH diagram ---
\centering
\safeimg[width=0.88\textwidth,height=4.8cm,keepaspectratio]{ch06-attention-mechanism.pdf}

\vspace{0.05cm}
\mlsysalert{Key:}{Doubling sequence length \textbf{quadruples} attention memory --- $O(N^2)$ is permanent.}

\end{frame}

\mlsysfocus{The Attention Trade-off}{%
\begin{minipage}{0.8\textwidth}\centering
\large
RNN: $O(1)$ memory, $O(N)$ sequential steps\\[0.5cm]
Transformer: $O(N^2)$ memory, $O(1)$ information depth\\[0.5cm]
\normalsize\textcolor{white}{Choose your constraint: time or space.}
\end{minipage}%
}

\begin{frame}{The Transformer Architecture}
\note{
% -- LINK: What prior concept connects to this slide
The focus frame just presented the RNN-vs-Transformer trade-off. This
slide dissects the Transformer block to show where compute actually goes.

% -- NARRATE: What to SAY while showing this slide
Point to the diagram: ``Two sub-layers per block. Multi-head attention
is 33\% of FLOPs but owns the O(N\^{}2) memory. The feed-forward MLP
is 67\% of FLOPs---it is the computational workhorse, not attention.''
Point to skip connections: ``These gradient highways enable stacking
100+ layers. Without them, training collapses beyond 20 layers.''
Point to parallelization card: ``All N tokens processed simultaneously---
this is why Transformers replaced RNNs for training throughput.''

% -- ENGAGE: Specific question, prediction, or task for THIS slide
Ask: ``Which sub-layer would you optimize first to speed up a Transformer:
attention or FFN?'' Give 15 seconds.
[Expected answer: FFN---it accounts for 67\% of FLOPs. But attention
dominates memory, so the answer depends on whether you are compute-bound
or memory-bound.]

% -- WARN: What students will get wrong on THIS topic
Common error: students assume attention is the expensive part because
it gets all the research attention. The FFN is actually the larger
compute cost per block.
Correct framing: attention dominates \emph{memory}; FFN dominates \emph{compute}.

% -- FLEX: [CORE] Students must understand the two-sub-layer anatomy.
[CORE] The 33/67 split recurs in Ch8 (training optimization).
IF SHORT: Skip the parallelization card; keep the diagram walkthrough.
}

\small
\begin{columns}[T]
  \begin{column}{0.48\textwidth}
    \safeimg[width=\textwidth,height=5.5cm,keepaspectratio]{ch06-transformer-block.pdf}
  \end{column}
  \begin{column}{0.48\textwidth}
    \textbf{Two sub-layers per block:}

    \vspace{0.15cm}
    \textcolor{computestroke}{\textbf{Multi-Head Attention}}\\
    {\footnotesize $\sim$33\% of FLOPs, $O(N^2)$ memory}

    \vspace{0.1cm}
    \textcolor{routingstroke}{\textbf{Feed-Forward MLP (4$\times$ hidden)}}\\
    {\footnotesize $\sim$67\% of FLOPs, dense MatMul}

    \vspace{0.1cm}
    \textcolor{datastroke}{\textbf{Skip + LayerNorm}}\\
    {\footnotesize Enables training 100+ layer networks}

    \vspace{0.15cm}
    \begin{mlsyscard}{computestroke}
    {\footnotesize \textbf{Parallelization:} All $N$ tokens processed simultaneously --- vs.\ RNN's $N$ sequential steps.}
    \end{mlsyscard}
  \end{column}
\end{columns}

\end{frame}

% --- ACTIVE LEARNING 2: Exercise ---
\begin{frame}{Your Turn: Attention Memory Budget}
\note{
% -- LINK: What prior concept connects to this slide
Students learned O(N\^{}2) attention memory two slides ago. Now they
must calculate the concrete memory cost for a real model configuration.

% -- NARRATE: What to SAY while showing this slide
``Time for math. A Transformer with 4096 tokens, 32 heads, 32 layers.
How much memory just for attention matrices?'' Present the problem (30s).
After 90 seconds, use peer instruction: if 30--70\% get it right, have
pairs discuss for 90s, then re-vote. Then reveal.

% -- ENGAGE: Specific question, prediction, or task for THIS slide
Students calculate individually (90s), then compare with a neighbor.
Per head: 4096\^{}2 * 2 bytes = 32 MB.
Total: 32 heads * 32 layers * 32 MB = 32 GB.
This exceeds a single GPU's memory for the attention matrix alone.

% -- WARN: What students will get wrong on THIS topic
Common error: students forget to multiply by number of heads AND layers.
They compute per-head correctly but miss the 32*32 multiplier.
IF STUCK: ``Start with one head, one layer. Then scale up.''

% -- FLEX: [CORE] This exercise makes O(N\^{}2) viscerally concrete.
[CORE] The 32 GB number motivates FlashAttention in Ch8.
IF AHEAD: ``What if we use FP32 instead of FP16?'' [Answer: 64 GB.]
IF SHORT: Show the solution immediately and walk through the arithmetic.
}

\footnotesize
\begin{columns}[T]
  \begin{column}{0.58\textwidth}
    {\small\bfseries Calculate Attention Memory}

    \vspace{0.1cm}
    A Transformer processes $N = 4{,}096$ tokens.
    \begin{itemize}\setlength\itemsep{0pt}
      \item Attention matrix: $N \times N$ in FP16 (2 bytes/entry)
      \item Model: 32 layers, 32 heads per layer
    \end{itemize}

    \vspace{0.1cm}
    \textbf{Calculate:} (1) Memory per attention matrix, (2) Total across all heads/layers

    {\footnotesize\textcolor{midgray}{(90 seconds --- compare with a neighbor)}}
  \end{column}
  \begin{column}{0.38\textwidth}
    \pause
    \begin{mlsyscard}{errorstroke}
    \textbf{Solution:}\\[0.05cm]
    {\footnotesize
    Per head: $4096^2 \times 2 = $ \textbf{32 MB}\\[0.05cm]
    Total: $32 \times 32 \times 32$ MB = \textbf{32 GB}\\[0.05cm]
    \textcolor{errorstroke}{\textbf{Exceeds a single GPU!}}
    }
    \end{mlsyscard}
  \end{column}
\end{columns}

\end{frame}

\begin{frame}{Transformer Compute Cost}
\note{
% -- LINK: Students just saw the Transformer block anatomy (33\% attention,
% 67\% FFN). This slide derives the exact FLOP count so students can
% calculate training and inference costs for any Transformer configuration.

% -- NARRATE: Walk through the formula.
``The total FLOPs for a Transformer forward pass scale as $12 \times L
\times H^2 \times S$ where $L$ is layers, $H$ is hidden dimension, and
$S$ is sequence length. For GPT-2 (12 layers, 768 hidden, 1024 tokens):
$\sim$7 GFLOPs per forward pass. GPT-3 (96 layers, 12288 hidden, 2048
tokens): $\sim$314 TFLOPs per forward pass. That is a 45,000$\times$
gap from GPT-2 to GPT-3.''

% -- ENGAGE: ``GPT-3 has 100$\times$ more parameters than GPT-2.
% Why is the FLOP count 45,000$\times$ more, not 100$\times$?''
% Give 15 seconds. [Expected answer: FLOPs scale with $H^2$, so doubling
% hidden dimension quadruples FLOPs. Plus longer sequences.]

% -- WARN: Students will confuse parameter count with FLOP count.
% Parameters scale as $O(H^2)$, FLOPs scale as $O(H^2 \times S)$.
% Sequence length is the multiplicative factor students forget.

% -- FLEX: [CORE] Quantitative FLOP analysis for Transformers.
% IF SHORT: Show the formula and the GPT-2 vs GPT-3 comparison only.
}

\footnotesize
\begin{columns}[T]
  \begin{column}{0.55\textwidth}
    \textbf{Transformer forward pass FLOPs:}

    \vspace{0.1cm}
    {\small $\text{FLOPs} \approx 12 \times L \times H^2 \times S$}

    \vspace{0.1cm}
    {\scriptsize
    \begin{tabular}{@{}ll@{}}
      $L$ & Number of layers \\
      $H$ & Hidden dimension \\
      $S$ & Sequence length \\
    \end{tabular}
    }

    \vspace{0.1cm}
    \textbf{Training} adds $\sim$2$\times$ (backward pass):
    $$\text{Train FLOPs} \approx 6 \times N \times D$$
    {\scriptsize where $N$ = parameters, $D$ = tokens processed.}
  \end{column}
  \begin{column}{0.42\textwidth}
    \renewcommand{\arraystretch}{1.1}
    {\scriptsize
    \begin{tabular}{@{}lrr@{}}
      \toprule
      & \textbf{GPT-2} & \textbf{GPT-3} \\
      \midrule
      Layers & 12 & 96 \\
      Hidden & 768 & 12,288 \\
      Seq len & 1,024 & 2,048 \\
      Params & 117M & 175B \\
      Fwd FLOPs & $\sim$7G & $\sim$314T \\
      \textbf{Ratio} & \multicolumn{2}{c}{\textbf{45,000$\times$}} \\
      \bottomrule
    \end{tabular}
    }

    \vspace{0.1cm}
    \begin{mlsyscard}{errorstroke}
    {\scriptsize FLOPs scale with $H^2 \times S$. Doubling hidden dimension \textbf{quadruples} compute.}
    \end{mlsyscard}
  \end{column}
\end{columns}

\end{frame}

% =============================================================================
\section{Sparse \& RecSys}
% =============================================================================

\begin{frame}{DLRM: Sparse Architectures}
\note{
% -- LINK: What prior concept connects to this slide
All previous architectures (MLP, CNN, RNN, Transformer) use dense operations.
DLRM is fundamentally different: the bottleneck is raw memory \emph{capacity},
not compute or bandwidth.

% -- NARRATE: What to SAY while showing this slide
``DLRM powers recommendation systems at Meta, Google, TikTok. The core
operation is not matrix multiplication---it is embedding table lookups.
Each user ID and item ID maps to a row in a table with 25 billion+ entries.
Each lookup is a random memory access, not a dense computation. This makes
DLRM the worst case for GPU caches---random access patterns defeat
spatial locality.''
Point to the table: ``Arithmetic intensity far below 1. The GPU's compute
units sit idle while memory fetches rows from terabytes of DRAM.''

% -- ENGAGE: Specific question, prediction, or task for THIS slide
Ask: ``Why can't you just put DLRM embedding tables in GPU HBM?''
Give 10 seconds. [Expected answer: tables are 100+ GB in production;
GPU HBM is 80 GB max. The tables do not fit.]

% -- WARN: What students will get wrong on THIS topic
Common error: students try to apply GPU optimization strategies
(batching, fusion) to DLRM. These do not help because the bottleneck
is capacity, not compute or bandwidth.
Correct framing: DLRM requires a fundamentally different system design---
distributed memory across many machines, SSD-backed embedding tables.

% -- FLEX: [CORE] DLRM represents a completely different bottleneck regime.
[CORE] Completes the five-lighthouse set.
IF SHORT: Cover the capacity bottleneck; skip the random access detail.
}

\footnotesize
\begin{columns}[T]
  \begin{column}{0.55\textwidth}
    \textbf{Deep Learning Recommendation Model}

    \vspace{0.05cm}
    \begin{itemize}\setlength\itemsep{0pt}
      \item \textbf{Data}: Categorical features (user ID, item ID)
      \item \textbf{Core op}: Embedding table lookups, not MatMul
      \item \textbf{Tables}: 25B+ entries, 100+ GB in production
      \item \textbf{Access}: Random, sparse --- worst case for caches
    \end{itemize}

    \vspace{0.1cm}
    \begin{mlsyscard}{errorstroke}
    Neither compute nor bandwidth but \alert{raw memory capacity} is the binding constraint.
    \end{mlsyscard}
  \end{column}
  \begin{column}{0.42\textwidth}
    \renewcommand{\arraystretch}{1.1}
    {\scriptsize
    \begin{tabular}{@{}lr@{}}
      \toprule
      \textbf{Property} & \textbf{Value} \\
      \midrule
      Embed.\ entries & 25B+ \\
      Table size & 100+ GB \\
      Arith.\ intensity & $\ll 1$ \\
      Access pattern & Random \\
      GPU utilization & Very low \\
      \bottomrule
    \end{tabular}
    }
  \end{column}
\end{columns}

\end{frame}

\begin{frame}{DLRM at Meta: Production Numbers}
\note{
% -- LINK: Students just learned DLRM's architectural properties. This slide
% grounds them with real production numbers from Meta's recommendation system.

% -- NARRATE: Walk through the numbers.
``Meta's DLRM serves billions of recommendations per day. The embedding
tables total 10+ TB across the fleet. Each inference touches 50--100 tables.
A single recommendation requires reading scattered rows from TB-scale
tables---the antithesis of dense MatMul.''

% -- ENGAGE: ``If each recommendation reads 100 embedding rows of 128
% floats each, how much data is accessed per inference?''
% Give 10 seconds. [Expected answer: 100 $\times$ 128 $\times$ 4 = 50 KB.
% But scattered across TB of memory, defeating all caching.]

% -- WARN: Students will try to apply GPU optimization strategies (batching,
% fusion) to DLRM. These help the dense MLP part but not the embedding
% lookups---the dominant cost.

% -- FLEX: [OPTIONAL] Concrete production numbers for DLRM.
% IF SHORT: Show the table for 20 seconds; emphasize the TB-scale challenge.
}

\footnotesize
\begin{columns}[T]
  \begin{column}{0.55\textwidth}
    \textbf{Meta DLRM Production System:}
    \begin{itemize}\setlength\itemsep{0pt}
      \item {\scriptsize Serves \textbf{billions} of recommendations/day}
      \item {\scriptsize Embedding tables: \textbf{10+ TB} across fleet}
      \item {\scriptsize Each inference: 50--100 table lookups}
      \item {\scriptsize DLRM accounts for \textbf{$>$50\%} of Meta's AI training compute}
    \end{itemize}

    \vspace{0.1cm}
    \begin{mlsyscard}{routingstroke}
    {\scriptsize \textbf{System design:} Distributed across 100s of machines. SSD-backed embedding tables. Custom hardware (ZionEX) with high-bandwidth memory pooling.}
    \end{mlsyscard}
  \end{column}
  \begin{column}{0.42\textwidth}
    \renewcommand{\arraystretch}{1.05}
    {\scriptsize
    \begin{tabular}{@{}lr@{}}
      \toprule
      \textbf{Metric} & \textbf{Value} \\
      \midrule
      Daily inferences & Billions \\
      Total embed.\ size & 10+ TB \\
      Tables per query & 50--100 \\
      Latency budget & $<$50 ms \\
      \% of AI compute & $>$50\% \\
      \bottomrule
    \end{tabular}
    }

    \vspace{0.1cm}
    \mlsysalert{Key:}{Recommendation is the largest AI workload at Meta---not LLMs.}
  \end{column}
\end{columns}
\mlsyscite{Naumov et al., ``Deep Learning Recommendation Model,'' 2019}

\end{frame}

% =============================================================================
\section{Building Blocks}
% =============================================================================

\begin{frame}{Shared Building Blocks Across Architectures}
\note{
% -- LINK: What prior concept connects to this slide
We have seen five architectures with different biases and bottlenecks.
Now we identify three building blocks that transfer across all of them---
the common infrastructure underneath the diversity.

% -- NARRATE: What to SAY while showing this slide
Point to each row: ``Skip connections: born in ResNets, now in every
Transformer block. They create gradient highways that enable training
beyond 20 layers. Normalization: rescales activations to prevent
explosion or collapse. Gating: learned switches that selectively retain
or forget state---LSTM invented this, Transformers inherited it.''
``These are prerequisites, not optimizations. Without skip connections,
you cannot train deep networks at all.''

% -- ENGAGE: Specific question, prediction, or task for THIS slide
Ask: ``What would happen to a 100-layer Transformer if you removed all
skip connections?'' Give 10 seconds.
[Expected answer: training would fail---gradients would vanish or explode
without the identity shortcut.]

% -- WARN: What students will get wrong on THIS topic
Common error: students treat these as optional performance tricks.
Correct framing: skip connections and normalization are structural
requirements for depth. Removing them breaks training, not just slows it.

% -- FLEX: [CORE] These three primitives appear in every modern architecture.
[CORE] Skip connections are referenced in the Key Takeaways.
IF SHORT: Focus on skip connections only; mention normalization and gating
in passing.
}

\footnotesize
\textbf{Three primitives that enable deep networks:}

\vspace{0.1cm}
\renewcommand{\arraystretch}{1.1}
{\scriptsize
\begin{tabular}{@{}lp{3cm}p{3cm}l@{}}
  \toprule
  \textbf{Block} & \textbf{What It Does} & \textbf{Why It Matters} & \textbf{Used In} \\
  \midrule
  \textbf{Skip Conn.} & $y = F(x) + x$ & Gradient highway, trains 100+ layers & ResNet, Transf. \\
  \textbf{Norm.} & Rescale activations & Prevents explosion/collapse & All modern nets \\
  \textbf{Gating} & Learned switches & Selectively retains/forgets state & LSTM, GRU, Transf. \\
  \bottomrule
\end{tabular}
}

\vspace{0.1cm}
\begin{mlsyscard}{datastroke}
These are \textbf{prerequisites}, not optimizations. Without skip connections, networks cannot train beyond $\sim$20 layers. Born in CNNs/RNNs, they transfer to every modern architecture.
\end{mlsyscard}

\end{frame}

\begin{frame}{Computational Primitives}
\note{
% -- LINK: What prior concept connects to this slide
Building blocks are the architectural patterns. Computational primitives
are the low-level operations those blocks reduce to on hardware.

% -- NARRATE: What to SAY while showing this slide
``Every architecture, no matter how complex, reduces to five types of
primitive operations. Dense MatMul---matrix times matrix---maps perfectly
to GPUs at 80--95\% utilization. MatVec is bandwidth-bound because one
operand is a vector. Sparse lookups (DLRM) are capacity-bound. Reductions
(softmax, LayerNorm) are bandwidth-bound because they read all data but
produce a small result.''
Point to the mapping rule: ``Maximize time spent in dense MatMul. That is
the engineering goal.''

% -- ENGAGE: Specific question, prediction, or task for THIS slide
Ask: ``Which primitive dominates during GPT-2 inference: MatMul or MatVec?''
Give 10 seconds. [Expected answer: MatVec---during autoregressive
generation, each step multiplies the weight matrix by a single token vector.]

% -- WARN: What students will get wrong on THIS topic
Common error: students assume all matrix operations are equally fast.
Correct framing: MatMul (matrix*matrix) achieves 80--95\% utilization;
MatVec (matrix*vector) achieves 10--30\%. The shape of the operands
determines GPU efficiency.

% -- FLEX: [OPTIONAL] Reinforces the roofline model concepts.
[OPTIONAL] Safe to abbreviate if running short.
IF SHORT: Show the table for 30 seconds; state the mapping rule; move on.
}

\scriptsize
\renewcommand{\arraystretch}{1.05}
\begin{tabular}{@{}lllrl@{}}
  \toprule
  \textbf{Primitive} & \textbf{Operation} & \textbf{Bottleneck} & \textbf{FLOPs/B} & \textbf{GPU Util.} \\
  \midrule
  \textbf{Dense MatMul} & Matrix $\times$ Matrix & Compute & 200+ & 80--95\% \\
  \textbf{MatVec} & Matrix $\times$ Vector & Bandwidth & $\sim$2 & 10--30\% \\
  \textbf{Sparse Lookup} & Embedding gather & Capacity & $\ll 1$ & $<$10\% \\
  \textbf{Reduction} & Softmax, LayerNorm & Bandwidth & 5--8 & 15--30\% \\
  \textbf{Elementwise} & ReLU, Add & Bandwidth & 0.1 & $<$5\% \\
  \bottomrule
\end{tabular}

\vspace{0.1cm}
\footnotesize
\mlsysconcept{Mapping rule:}{Dense MatMul achieves 80--95\% GPU utilization. Everything else: 10--50\%. Maximize time in dense MatMul.}

\end{frame}

% =============================================================================
\section{Selection}
% =============================================================================

\begin{frame}{Architecture Selection Framework}
\note{
% -- LINK: What prior concept connects to this slide
Students now know five architectures, their biases, their bottlenecks,
and their primitives. This slide provides the decision framework for
choosing among them.

% -- NARRATE: What to SAY while showing this slide
Walk through the flowchart: ``Two phases. Phase 1: match your data type
to architecture candidates---spatial data suggests CNN, sequential suggests
RNN or Transformer. Phase 2: validate against deployment constraints---
memory, latency, power. Failing ANY constraint means the model has zero
utility regardless of accuracy. The retry loop forces scaling down or
switching architecture entirely.''

% -- ENGAGE: Specific question, prediction, or task for THIS slide
Ask: ``A model achieves 99\% accuracy but exceeds the memory budget by
10\%. Is it deployable?'' Give 10 seconds.
[Expected answer: no---violating any deployment constraint means zero
utility. 99\% accuracy is irrelevant if it cannot run.]

% -- WARN: What students will get wrong on THIS topic
Common error: students evaluate architecture by accuracy first, then
try to fit it into constraints afterward.
Correct framing: constraints are hard gates. Filter by constraints first,
then maximize accuracy within the feasible set.

% -- FLEX: [CORE] This flowchart is the practical tool students take away.
[CORE] Directly used in the Discussion exercise next.
IF AHEAD: ``What if no architecture meets all constraints?'' [Answer:
change the constraints or change the problem.]
}

% --- Layout: FULL-WIDTH diagram ---
\centering
\safeimg[width=0.92\textwidth,height=4.8cm,keepaspectratio]{ch06-architecture-selection-flowchart.pdf}

\vspace{0.1cm}
\mlsysalert{Rule:}{A model that violates any deployment constraint has \textbf{zero utility} --- regardless of accuracy.}

\end{frame}

\begin{frame}{Architecture Complexity Comparison}
\note{
% -- LINK: What prior concept connects to this slide
The selection flowchart identifies candidates. This comparison table
provides the quantitative data to choose among them.

% -- NARRATE: What to SAY while showing this slide
``No single best architecture. CNNs dominate spatial perception at 80--95\%
GPU utilization. Transformers master reasoning but cost O(N\^{}2) memory.
RNNs have zero parallelism---30--50\% utilization. DLRM needs terabytes
of memory for embedding tables.''
Point to the bottom card: ``Match bias to data, validate against constraints.
That is the full decision procedure.''

% -- ENGAGE: Specific question, prediction, or task for THIS slide
Ask: ``If you have spatial data and an edge device with 4 GB RAM, which
column do you pick?'' Give 10 seconds.
[Expected answer: CNN---spatial bias matches the data, and CNNs fit in
small memory budgets.]

% -- WARN: What students will get wrong on THIS topic
Common error: students default to Transformers for everything because
of recent hype.
Correct framing: Transformers are the most flexible but the most expensive.
For spatial data on edge, a CNN or MobileNet variant is the right choice.

% -- FLEX: [OPTIONAL] The table summarizes what students already know.
[OPTIONAL] If running short, let students read it; do not narrate every cell.
IF SHORT: Point to the CNN and Transformer columns only; move to Discussion.
}

\scriptsize
\renewcommand{\arraystretch}{1.05}
\begin{tabular}{@{}llllll@{}}
  \toprule
  & \textbf{MLP} & \textbf{CNN} & \textbf{RNN} & \textbf{Transformer} & \textbf{DLRM} \\
  \midrule
  \textbf{Bias}    & None    & Spatial      & Temporal     & Learned      & Categorical \\
  \textbf{Params}  & $O(d^2)$  & $O(K^2 C)$ & $O(d^2)$   & $O(d^2)$   & $O(\text{vocab})$ \\
  \textbf{Memory}  & Linear  & Linear       & $O(1)$/step  & $O(N^2)$     & TB-scale \\
  \textbf{Parallel}& Full    & Full         & None         & Full         & Full \\
  \textbf{GPU util.}& 80--90\% & 80--95\%   & 30--50\%     & 50--80\%     & 10--30\% \\
  \bottomrule
\end{tabular}

\vspace{0.1cm}
\footnotesize
\begin{mlsyscard}{crimson}
No single ``best'' architecture. CNNs dominate spatial perception. Transformers master reasoning at quadratic cost. DLRM needs TB-scale memory. \textbf{Match bias to data, validate against constraints.}
\end{mlsyscard}

\end{frame}

% --- Quantitative Exercise ---
\begin{frame}{Your Turn: Jetson Orin Throughput}
\note{
% -- LINK: Students just saw the selection framework and comparison table.
% This exercise forces them to apply quantitative reasoning to a real
% deployment constraint.

% -- NARRATE: Present the problem (30s). Individual work (60s).
``You need to process 1000 images per second on a Jetson Orin NX (100 TOPS
INT8). ResNet-50 takes 5 ms per image. ViT-B takes 12 ms. Which fits?
What batch size is needed?''

% -- ENGAGE: Students calculate individually (60s), then compare.
% ResNet-50: 1/5ms = 200 img/s per stream. Need 1000/200 = 5 parallel
% streams. At batch=5: 5$\times$5ms = 25ms per batch of 5 = 200 batches/s
% = 1000 img/s. Fits!
% ViT-B: 1/12ms = 83 img/s. Need 1000/83 = 12 streams. But 12$\times$340MB
% = 4 GB weights alone. 32 GB RAM, activations scale too. Tight.

% -- WARN: Students will divide peak TOPS by model FLOPs.
% Correct framing: use measured latency, not theoretical peak.
% Real utilization on edge is 30--50\% of peak TOPS.

% -- FLEX: [CORE] Practical deployment calculation.
% IF SHORT: Show the solution immediately; emphasize the memory constraint.
}

\footnotesize
\begin{columns}[T]
  \begin{column}{0.55\textwidth}
    {\small\bfseries Edge Deployment Calculation}

    \vspace{0.1cm}
    \textbf{Target}: 1,000 images/sec on \textbf{Jetson Orin NX}\\
    (100 TOPS INT8, 32 GB LPDDR5)

    \vspace{0.1cm}
    Measured inference latency (batch=1):
    \begin{itemize}\setlength\itemsep{0pt}
      \item ResNet-50: \textbf{5 ms/image}
      \item ViT-B: \textbf{12 ms/image}
    \end{itemize}

    \vspace{0.1cm}
    \textbf{Calculate:} Which model meets the target?\\
    What batch size is needed?

    {\footnotesize\textcolor{midgray}{(60 seconds --- then compare)}}
  \end{column}
  \begin{column}{0.42\textwidth}
    \pause
    \begin{mlsyscard}{datastroke}
    \textbf{Solution:}\\[0.05cm]
    {\scriptsize
    ResNet: $1/5\text{ms} = 200$ img/s\\
    Need batch $\geq 5$ for 1000 img/s\\
    Memory: $5 \times 98$ MB = 490 MB\\
    \textcolor{datastroke}{\textbf{Fits easily!}}\\[0.1cm]
    ViT-B: $1/12\text{ms} = 83$ img/s\\
    Need batch $\geq 12$ for 1000 img/s\\
    Memory: $12 \times 340$ MB = 4 GB\\
    \textcolor{routingstroke}{\textbf{Tight on 32 GB}}
    }
    \end{mlsyscard}
  \end{column}
\end{columns}

\end{frame}

% --- ACTIVE LEARNING 3: Discussion ---
\begin{frame}{Discussion: Pick the Architecture}
\note{
% -- LINK: What prior concept connects to this slide
Students just saw the selection framework and comparison table. This
exercise forces them to apply both to a concrete scenario with real
constraints.

% -- NARRATE: What to SAY while showing this slide
``Turn and talk with your neighbor for 90 seconds. A startup has a
vision model on an edge device: 4 GB RAM, 100 ms latency. They trained
a ViT with 86M parameters requiring 340 MB weights. Should they deploy
it? If not, what architecture should they use?''
After 90 seconds, take a show-of-hands poll across the four options.

% -- ENGAGE: Specific question, prediction, or task for THIS slide
Students discuss in pairs for 90 seconds. The question is falsifiable:
ViT attention memory at inference will exceed the budget for long sequences,
and the 100 ms latency target is hard to meet without compilation.
Cold-call 2--3 pairs for their reasoning.
[Expected answer: ViT is risky---quadratic memory on 4 GB is tight.
CNN or MobileNet variant is the safer choice for spatial data on edge.]

% -- WARN: What students will get wrong on THIS topic
Common error: students say ``just quantize the ViT to INT8.''
Correct framing: quantization halves weight memory (170 MB) but does not
fix the O(N\^{}2) attention memory scaling. At high resolution, attention
memory can exceed the budget even with quantized weights.

% -- FLEX: [CORE] This is the chapter's primary application exercise.
[CORE] Must complete this discussion before moving to Wrap-Up.
IF SHORT: Do a show-of-hands poll (30s) instead of full turn-and-talk.
}

\centering
\vspace{0.5cm}
{\large\bfseries Turn and Talk \textcolor{midgray}{(90 seconds)}}

\vspace{0.4cm}
{\normalsize A startup deploys a vision model on an edge device\\
with \textbf{4 GB RAM} and a \textbf{100 ms latency} target.\\[0.2cm]
Their ML team trained a Vision Transformer (ViT)\\
with 86M parameters requiring 340 MB weights.}

\vspace{0.3cm}
\alert{Should they deploy this model? If not, what architecture?}

\vspace{0.2cm}
\begin{columns}[c]
  \begin{column}{0.20\textwidth}\centering\footnotesize Deploy ViT\end{column}
  \begin{column}{0.20\textwidth}\centering\footnotesize Switch to CNN\end{column}
  \begin{column}{0.20\textwidth}\centering\footnotesize Switch to MobileNet\end{column}
  \begin{column}{0.20\textwidth}\centering\footnotesize Quantize ViT\end{column}
\end{columns}

\end{frame}

% =============================================================================
\section{Wrap-Up}
% =============================================================================

\begin{frame}{Fallacies}
\note{
% -- LINK: What prior concept connects to this slide
Students applied the selection framework in the discussion exercise.
These fallacies show what goes wrong when people skip that framework.

% -- NARRATE: What to SAY while showing this slide
Read each fallacy and its rebuttal. Emphasize the numbers:
``CNN hits 99\% on MNIST with 420K params vs.\ MLP's 20M---complexity
is not a proxy for performance.''
``Transformer at 50 ms on A100 becomes 2,000 ms on mobile---a 40x
slowdown because mobile lacks HBM.''
``MobileNet uses 14x fewer FLOPs but runs slower on datacenter GPUs---
FLOPs do not predict speed.''
``KV cache at 2048 tokens, 32 layers, 32 heads: 537 MB per user.
Four concurrent users rival model weight memory.''

% -- FLEX: [CORE] Fallacies reinforce the chapter's quantitative reasoning.
[CORE] Each fallacy maps to a specific concept from the lecture.
IF SHORT: Cover fallacies 1 and 3 (complexity and FLOPs); skip 2 and 4.
}

\footnotesize
\textbf{Fallacy:} \textit{More complex architectures always perform better.}\\
CNN: 99\% on MNIST with 420K params. MLP: 98\% with 20M params. Defaulting to Transformers for tabular data wastes 5--10$\times$ resources.

\vspace{0.08cm}
\textbf{Fallacy:} \textit{Architecture performance transfers across hardware.}\\
Transformer at 50 ms on A100 $\to$ 2,000 ms on mobile SoC (40$\times$ slowdown, no HBM).

\vspace{0.08cm}
\textbf{Fallacy:} \textit{FLOPs determine inference speed.}\\
MobileNet: 14$\times$ fewer FLOPs than ResNet but runs \emph{slower} on datacenter GPUs (low intensity starves ALUs).

\vspace{0.08cm}
\textbf{Fallacy:} \textit{KV cache is a minor serving cost.}\\
Seq 2,048, 32 layers $\times$ 32 heads: each user needs $\sim$537 MB KV cache. At 4 users, rivals model weights.

\end{frame}

\begin{frame}{Pitfalls}
\note{
% -- LINK: What prior concept connects to this slide
Fallacies are wrong beliefs; pitfalls are operational mistakes that follow
from correct knowledge applied carelessly.

% -- NARRATE: What to SAY while showing this slide
``Pitfall 1: selecting by accuracy alone. A Transformer at sequence 2048
needs 16x more memory than at 512. You hit latency SLAs, you drop to
25\% GPU efficiency.''
``Pitfall 2: mixing architectural patterns without profiling. Adding
attention to a CNN flushes feature maps from cache---ResNet drops from
250 to 80 img/s, a 3x hit.''
``Pitfall 3: designing for training hardware. Your 8xA100 cluster has
640 GB; your edge device has 4 GB---a 160x gap that quantization alone
cannot bridge.''

% -- FLEX: [CORE] Operational pitfalls complement the conceptual fallacies.
[CORE] These are real-world mistakes students will encounter in projects.
IF SHORT: Cover pitfalls 1 and 3; skip pitfall 2.
}

\footnotesize
\textbf{Pitfall:} \textit{Selecting architectures based solely on accuracy.}\\
Transformers: seq 2,048 needs 16$\times$ more memory than 512. Miss latency SLAs (100 ms $\to$ 500 ms), 25\% GPU efficiency.

\vspace{0.08cm}
\textbf{Pitfall:} \textit{Combining patterns without profiling interactions.}\\
Attention + CNN flushes feature maps from cache. ResNet: 250 img/s $\to$ 80 img/s (3$\times$ drop).

\vspace{0.08cm}
\textbf{Pitfall:} \textit{Designing for training hardware, ignoring deployment.}\\
8$\times$ A100 cluster (640 GB) $\to$ 4 GB edge device = 160$\times$ gap. Requires architectural changes, not just quantization.

\end{frame}


\begin{frame}{Muddiest Point}
\note{
% -- NARRATE: What to SAY while showing this slide
``One-minute paper. Write down the muddiest point---the concept that was
most confusing or unclear. Hand it in on your way out or submit digitally.
I will open next lecture by addressing the top 2--3 muddiest points.''

% -- FLEX: [CORE] Diagnostic tool for the instructor.
[CORE] Collect responses to open the next lecture with targeted clarification.
IF SHORT: Ask students to submit digitally after class instead of in person.
}

\centering
\Large\bfseries One-minute paper\\[0.5cm]
\normalsize Write down the \textbf{muddiest point} from today's lecture.\\[0.2cm]
{\small What concept was most confusing or unclear?}\\[0.5cm]
{\footnotesize\textcolor{midgray}{(Hand in on your way out --- or submit digitally)}}
\end{frame}

% --- RETRIEVAL PRACTICE ---
\begin{frame}{What Were the Key Ideas?}
\note{
% -- NARRATE: What to SAY while showing this slide
``Close your notes. Write down the four most important concepts from today.
90 seconds, no peeking.'' Walk around the room to monitor engagement.
Do NOT show the Key Takeaways slide yet.

% -- FLEX: [CORE] Retrieval practice is the highest-impact learning activity.
[CORE] Students must write before seeing the answer slide.
IF SHORT: Reduce to 60 seconds.
}

\centering
\vspace{1.5cm}
{\Large\bfseries Close your notes.}

\vspace{0.8cm}
{\large Write down the \textbf{4 most important concepts} from today.}

\vspace{0.8cm}
{\normalsize\textcolor{midgray}{90 seconds --- no peeking.}}

\end{frame}

\begin{frame}{Key Takeaways}
\note{
% -- LINK: What prior concept connects to this slide
Students just did retrieval practice. This slide reveals the answers
so they can compare against their own lists.

% -- NARRATE: What to SAY while showing this slide
Walk through each bullet with its quantitative anchor:
``Inductive bias unifies all architectures---locality, sequence, global context.''
``Arithmetic intensity spans 1000x from ResNet to GPT-2.''
``O(N\^{}2) attention is permanent---it constrains deployment.''
``Five lighthouses, five bottlenecks: compute, bandwidth, capacity, latency, power.''
``FLOPs do not equal speed---MobileNet proves it.''
``Skip connections enable depth---prerequisite, not optimization.''
``Architecture equals deployment---the model determines your infrastructure.''

% -- FLEX: [CORE] Consolidates the lecture.
[CORE] Read all bullets; emphasize the numbers.
IF SHORT: Read bullets 1, 3, and 7; let students read the rest.
}

\footnotesize
\begin{itemize}\setlength\itemsep{0pt}
  \item \textbf{Inductive bias unifies all architectures}: Locality (CNN), sequence (RNN), global context (Transformer). Bias trades generality for efficiency.
  \item \textbf{Arithmetic intensity determines bottleneck}: High (CNN) = compute-bound. Low (GPT-2) = memory-bound. 1000$\times$ gap.
  \item \textbf{$O(N^2)$ is permanent}: Transformer attention scales quadratically. This constrains deployment.
  \item \textbf{Lighthouses isolate bottlenecks}: ResNet (compute), GPT-2 (bandwidth), DLRM (capacity), MobileNet (latency), KWS (power).
  \item \textbf{FLOPs $\neq$ speed}: MobileNet uses 14$\times$ fewer FLOPs but runs slower on datacenter GPUs.
  \item \textbf{Skip connections enable depth}: Prerequisites for networks beyond $\sim$20 layers.
  \item \textbf{Architecture = deployment}: The model determines memory, latency, power, and cost.
\end{itemize}

\end{frame}

\begin{frame}{References}
\note{
% -- NARRATE: What to SAY while showing this slide
``These are the canonical papers behind today's content. He et al.\ for
ResNets, Vaswani et al.\ for Transformers, Howard et al.\ for MobileNets.
Read at least the abstracts before next class.''

% -- FLEX: [OPTIONAL] Reference slides are for student self-study.
[OPTIONAL] Do not read citations aloud; point students to the slide.
}

\small
\mlsysref{He+16}{K. He et al. ``Deep Residual Learning for Image Recognition.'' CVPR 2016.}
\mlsysref{Vaswani+17}{A. Vaswani et al. ``Attention Is All You Need.'' NeurIPS 2017.}
\mlsysref{Howard+17}{A. Howard et al. ``MobileNets.'' arXiv 2017.}
\mlsysref{Naumov+19}{M. Naumov et al. ``Deep Learning Recommendation Model (DLRM).'' arXiv 2019.}
\mlsysref{Radford+19}{A. Radford et al. ``Language Models are Unsupervised Multitask Learners.'' 2019.}
\mlsysref{Bahdanau+14}{D. Bahdanau et al. ``Neural Machine Translation by Jointly Learning to Align.'' 2014.}

\end{frame}

\begin{frame}{Next Lecture: ML Frameworks}
\note{
% -- LINK: What prior concept connects to this slide
We now have the architectural blueprints---five lighthouses, their biases,
their bottlenecks. The next chapter asks: how do we actually build and
run these architectures on hardware?

% -- NARRATE: What to SAY while showing this slide
``We have the blueprints. Now we need the tools. PyTorch, TensorFlow,
and JAX translate these high-level graphs into low-level kernels that
run on silicon. The contract with physics that we signed today? The
framework compiler is what enforces it.''

% -- FLEX: [CORE] Forward hook must connect to Ch7.
[CORE] Keep to 1 minute. End on the compiler metaphor.
}

\footnotesize
\centering
\vspace{0.3cm}
{\large We have the \textbf{blueprints}. Now we need the \textbf{tools}.}

\vspace{0.4cm}
\begin{columns}[c]
  \begin{column}{0.30\textwidth}
    \centering
    \begin{mlsyscard}{computestroke}
    {\normalsize\bfseries PyTorch}\\[0.1cm]
    {\scriptsize Eager execution\\Dynamic graphs}
    \end{mlsyscard}
  \end{column}
  \begin{column}{0.30\textwidth}
    \centering
    \begin{mlsyscard}{datastroke}
    {\normalsize\bfseries TensorFlow}\\[0.1cm]
    {\scriptsize Graph-based\\Production serving}
    \end{mlsyscard}
  \end{column}
  \begin{column}{0.30\textwidth}
    \centering
    \begin{mlsyscard}{routingstroke}
    {\normalsize\bfseries JAX}\\[0.1cm]
    {\scriptsize Functional transforms\\XLA compilation}
    \end{mlsyscard}
  \end{column}
\end{columns}

\vspace{0.3cm}
{\normalsize How do frameworks translate architecture graphs\\into the low-level kernels that run on silicon?}\\[0.1cm]
{\small\textbf{The ``contract with physics'' signed here is enforced by the compiler.}}

\end{frame}


% =============================================================================
% BACKUP SLIDES
% =============================================================================
\appendix

\begin{frame}{Backup: KV Cache Memory Calculation}
\note{
% -- NARRATE: Use if students want to understand KV cache scaling in detail.
Per layer, per head: $2 \times N \times d_k$ values (K and V).
With $N=4096$, $d_k=64$, FP16: $2 \times 4096 \times 64 \times 2 = 1$ MB per head per layer.
32 heads $\times$ 32 layers = 1 GB just for KV cache at 4K context.

% -- FLEX: [OPTIONAL] Backup slide for advanced questions.
}
\small
Use if students want to understand the KV cache scaling.\\small Per layer, per head: $2 \times N \times d_k$ values (K and V).\\ With $N=4096$, $d_k=64$, FP16: $2 \times 4096 \times 64 \times 2 = 1$ MB per head per layer.\\ 32 heads $\times$ 32 layers = \textbf{1 GB} just for KV cache at 4K context.
\end{frame}

\begin{frame}{Backup: Why FLOPs Do Not Predict Speed}
\note{
% -- NARRATE: Use if students question the MobileNet paradox.

% -- FLEX: [OPTIONAL] Backup slide for the FLOPs paradox.
}
\small
Use if students question the MobileNet paradox.\\
MobileNet: 300 MFLOPs but low arithmetic intensity ($\sim$21 FLOPs/byte).\\
ResNet-50: 4.1 GFLOPs but high intensity ($\sim$40 FLOPs/byte).\\
On datacenter GPUs, ResNet saturates tensor cores; MobileNet starves them.\\
Result: MobileNet can be \emph{slower} despite 14$\times$ fewer FLOPs.
\end{frame}

\begin{frame}{Backup: The Architecture Arms Race}
\note{
% -- NARRATE: War story about architecture selection gone wrong.
% Use when students ask ``does architecture choice really matter in practice?''

% -- FLEX: [OPTIONAL] Motivational backup for industry context.
}
\footnotesize
\textbf{War Story: When Architecture Meets Production}

\vspace{0.1cm}
In 2020, a major cloud provider migrated their image classification pipeline from ResNet-50 to ViT-B, expecting accuracy gains. Results:

\vspace{0.05cm}
\scriptsize
\renewcommand{\arraystretch}{1.05}
\begin{tabular}{@{}lrr@{}}
  \toprule
  & \textbf{ResNet-50} & \textbf{ViT-B} \\
  \midrule
  Accuracy & 76.1\% & 77.9\% (+1.8\%) \\
  Latency (A100) & 3.2 ms & 8.7 ms \\
  Memory/request & 98 MB & 340 MB \\
  Serving cost & \$1.00 & \$2.70 \\
  \bottomrule
\end{tabular}

\vspace{0.05cm}
\footnotesize
\begin{mlsyscard}{errorstroke}
{\scriptsize \textbf{Lesson}: 1.8\% accuracy gain cost 2.7$\times$ more to serve. The team reverted to ResNet-50 with better data augmentation, achieving 77.5\% accuracy at the original cost.}
\end{mlsyscard}
\end{frame}

\begin{frame}{Backup: Architecture Memory Cheat Sheet}
\note{
% -- NARRATE: Quick reference for architecture memory scaling.
% -- FLEX: [OPTIONAL] Reference slide for assignments and projects.
}
\scriptsize
\renewcommand{\arraystretch}{1.05}
\begin{tabular}{@{}lllll@{}}
  \toprule
  \textbf{Architecture} & \textbf{Weight Memory} & \textbf{Activation Memory} & \textbf{Total (batch=32)} & \textbf{Fits On} \\
  \midrule
  MLP (784-1024-10) & 3.2 MB & 128 KB $\times$ B & $\sim$7 MB & Any \\
  ResNet-50 & 98 MB & 2 MB $\times$ B & $\sim$162 MB & Mobile \\
  GPT-2 (117M) & 468 MB & 50 MB $\times$ B & $\sim$2 GB & Edge GPU \\
  GPT-2 XL (1.5B) & 6 GB & 200 MB $\times$ B & $\sim$12 GB & Datacenter \\
  DLRM (production) & 100+ GB & Minimal & 100+ GB & Fleet \\
  \bottomrule
\end{tabular}

\vspace{0.1cm}
\begin{mlsyscard}{crimson}
{\scriptsize \textbf{Rule of thumb}: Weights determine if the model \emph{fits}. Activations determine if a \emph{batch} fits. Architecture determines both.}
\end{mlsyscard}
\end{frame}

\end{document}
